\documentclass[track=junior,level=2]{kidscode}
\begin{document}
\kctitle{Notes for Families}{What we are doing, how we keep children safe, and how you can help at home}

\section{Dear families,}
Welcome to Game Makers! Once a week, for 60 minutes, your child will learn to make their own
computer games, in a small group of up to four children (about 8--9 years old, US Grades 2--3).
This level follows Little Robots, so your child already knows how to snap Scratch blocks together,
use repeat, start code with a click or a key, and fix a bug. If your child is new, please tell me:
the first session starts with a short, friendly \enquote{show me} warm-up so we know where to begin.

In ten weeks children turn Scratch scripts into games: a sprite that keeps moving, an apple that
falls into a bowl, a score that goes up, surprises that nobody can guess, and an ending that says
\enquote{You win!} or \enquote{Game over}.
\begin{itemize}[itemsep=2pt]
  \item \textbf{Weeks 1--4:} forever loops, \enquote{if touching}, keeping score, random surprises.
  \item \textbf{Week 5:} a play day.
  \item \textbf{Weeks 6--8:} messages between sprites, making our own blocks, a start, a win and a game over.
  \item \textbf{Weeks 9--10:} the project, \enquote{My First Arcade Game}, planned on a game plan
        card and tested by another player, and an arcade showcase where your family plays the games
        and every child gets a certificate.
\end{itemize}
Every session has the same parts: a warm-up game or short video, a short demonstration
(\enquote{Show me}), building in pairs (\enquote{Make it}) with a wiggle break in the middle, and a
Share where children play each other's games. Children work as a pair: the \textbf{Driver} holds
the mouse and the \textbf{Helper} points and reads. A third child, or anyone testing a game, is the
\textbf{Player}. They swap jobs when the timer runs out.

\section{Our safety and privacy promises}
\begin{safety}[What we promise]
\begin{itemize}[itemsep=2pt]
  \item \textbf{Offline.} We use the Scratch app on the computer, not the website. Nothing needs the internet.
  \item \textbf{No accounts.} Nobody signs up, logs in or gives an email address.
  \item \textbf{Nicknames only.} Games show a nickname or a first name, never a full name,
        address, school, age or photo.
  \item \textbf{Nothing is uploaded.} Games stay on the computer. The arcade showcase is at home,
        with family only, on the screen.
  \item \textbf{Fair and kind games.} We talk about fair rules, kind feedback, being a good sport and
        games that do not trick players into \enquote{just one more go}.
  \item \textbf{Healthy screen time.} Every session has a wiggle break, and children learn to stop
        when the timer says so, even in the middle of a game.
  \item \textbf{For everyone.} All pages use color-blind-safe colors with words and icons, and every
        page comes in a dyslexia-friendly edition. Please ask.
\end{itemize}
\end{safety}

\section{What your child needs}
\begin{itemize}[itemsep=2pt]
  \item \textbf{A Mac with Scratch.} The free Scratch 3 app from Level 1. No account is needed.
  \item \textbf{A mouse and a keyboard.} Games are played with the arrow keys.
  \item \textbf{60 minutes, once a week,} at a time when your child is not too tired or hungry.
  \item \textbf{A grown-up nearby,} to read along and to be the Helper or the Player when there is one child.
\end{itemize}

\section{How to help at home}
You do not need to know anything about coding to help.
\begin{itemize}[itemsep=3pt]
  \item \textbf{Play the play-at-home card.} Each week there is an optional 5--10 minute game, mostly
        without a computer. Stop while it is still fun.
  \item \textbf{Be a Player.} When your child shows you a game, play it honestly and say what
        happened. Then give \enquote{two stars and a wish}: two things you liked and one idea.
  \item \textbf{Praise effort and good ideas.} \enquote{You found out why the score did not reset!}
        means more than \enquote{You won!}
  \item \textbf{Bugs are normal.} Ask \enquote{What did you want it to do? What did it do instead?}
  \item \textbf{Hands off the mouse.} If you help, be the Helper: point and use words, and let your child click.
  \item \textbf{Talk about stopping.} A forever loop never stops on its own, but we do. Agree on a
        stop time before your child starts playing any game.
\end{itemize}

\section{Ask your child about\dots}
One question for each week. They make great conversation starters at dinner.

\begin{center}\small\renewcommand{\arraystretch}{1.45}
\begin{tabularx}{\linewidth}{@{}>{\bfseries}p{1.0cm}>{\RaggedRight}p{4.2cm}>{\RaggedRight\arraybackslash}X@{}}
\toprule
Week & Topic & Ask your child\dots\\
\midrule
1 & Forever and Bounce & \enquote{What is the difference between repeat and forever? How do you stop a forever loop?}\\
2 & If Touching\dots & \enquote{What happens in your game when the apple touches the bowl?}\\
3 & Keep Score & \enquote{What is a variable? How does your game know the score?}\\
4 & Surprise! Random & \enquote{What did you make random? Why does that make a game more fun?}\\
5 & Play Day and Show Me & \enquote{What did you get better at? Which bug was the trickiest?}\\
6 & Messages & \enquote{How can one sprite tell the others to start?}\\
7 & My Own Blocks & \enquote{Which block did you make yourself? What is inside it?}\\
8 & Start, Win and Lose & \enquote{How do you win your game? How do you lose?}\\
9 & Project Studio & \enquote{What is your arcade game about? What did your Player find?}\\
10 & Arcade Showcase & \enquote{Why did you make your game this way? What would you make next?}\\
\bottomrule
\end{tabularx}
\end{center}

\section{How to celebrate}
\begin{itemize}[itemsep=3pt]
  \item \textbf{Come to the arcade showcase} in week 10. Your child will show the game, let you play
        it, tell you what one script does and why they made a choice. You get sticker feedback cards to fill in.
  \item \textbf{Certificates.} Every child who finishes Game Makers gets a printed certificate.
  \item \textbf{A fridge gallery.} Ask your child to draw their game plan card or a screen from
        their game for the fridge. Any photos or screenshots stay at home, not online.
  \item \textbf{Celebrate effort.} \enquote{You asked someone to test it and fixed what they found!}
        is the best thing a young game maker can hear.
\end{itemize}

\bigskip
Thank you for coding with us. If anything would help your child learn (for example bigger text,
more breaks or a different time of day), please tell me.

\bigskip
With warm wishes,\par\bigskip
\rule{6cm}{0.4pt}\par
{\small (Your coding club host)}
\end{document}
